\reviewsection{What Students Taught Me About Behavior, Communication, and Design}

The National Education Association guide became most useful when I connected it to moments from my classroom \citep{nea2020guide}. During technology, coding, and robotics activities, I saw how movement, repetition, avoidance, silence, leaving a group, or refusing a material could be interpreted as defiance when the actual issue was overload, uncertainty, lack of control, or inaccessible directions. I also saw how quickly behavior changed when the environment changed.

When I placed only the needed materials on the table, modeled the first step, reduced the number of verbal directions, clarified peer roles, or offered a quieter starting point, students often re-entered the activity without a lecture, consequence, or reward. Reading Rockets, the educator guide, and Denning and Moody all support this proactive move from waiting for failure to designing access in advance \citep{readingRocketsASD,iris2016asd,denningMoody2018}. Denning and Moody are especially helpful because they name a problem I saw repeatedly: classrooms are often organized around a presumed typical learner, and then individual students are treated as disruptions when they need the environment to change.

The inclusion sources validate the same point from another direction. \citet{bock2016inclusion} distinguish presence from meaningful participation and emphasize that teacher knowledge, attitudes, training, and administrative support affect whether inclusion becomes real. \citet{smith2011inclusion} likewise treats inclusion as planned, supported work rather than a hopeful placement decision. These sources strengthen my claim that inclusion requires structure and accountability. At the same time, their intervention-oriented language must be read alongside Conn and neurodiversity scholarship so that ``support'' does not quietly become training a student to appear more typical.

The strongest question I can ask is, \textit{``What might this student be communicating, and what barrier is present?''} That change in language changes the intervention. It moves me away from controlling the student and toward redesigning the conditions.

This also changed how I understand independence. Support should not mean doing the task for a student or making every decision on the student's behalf. The strongest scaffolds make the goal visible and the entry point accessible while preserving opportunities for choice, productive struggle, communication, revision, joy, and repair. A visual sequence, partner role, sensory option, communication board, model, or quiet reset should reveal competence rather than replace it.

\reviewsection{What Counts as a Good Outcome?}

My placement taught me to celebrate when students could enter an activity, communicate a need, recover from frustration, contribute an idea, or remain connected in their own way. Module 1 pushed me to ask whether schools sometimes define a ``good outcome'' too narrowly: task completion, quiet behavior, independence from support, employment, or the ability to resemble a typical peer.

The transition reading expands support beyond the immediate classroom. \citet{stuartCollins2019transition} connect visual supports, technology, self-determination, and transition planning to adult and community participation. Their work validates the importance of preparing for the future, but it also changes what preparation should mean. Future planning should begin with the student's goals, communication, preferences, relationships, and desired community life---not only with adult assumptions about productivity.

\textit{Best Kept Secret} makes that institutional problem visible through students nearing the end of public education \citep{buck2013best}. The film and transition reading together oppose a classroom-only understanding of access. A school has not fully supported agency if the student leaves with no meaningful pathway into adult services, relationships, community, continued learning, or self-directed life. Independence can be valuable, but interdependence is not failure. A good outcome is not a life without support; it is a life in which support does not erase authorship.

\reviewsection{Why Credibility, Culture, and Identity Cannot Be Separated}

My lived experience makes it impossible for me to think about developmental difference as an isolated variable. I am Latino, bilingual, neurodivergent, chronically ill, and shaped by trauma. Different institutions have interpreted the same traits differently depending on which part of me they noticed first. Sometimes persistence was seen as intensity. Sometimes difficulty with an inaccessible process was treated as lack of responsibility. Sometimes communication style mattered more than the content of what I was saying.

Students also move through the world with multiple identities. An autistic learner may be multilingual, racialized, gender-diverse, economically marginalized, physically disabled, or part of a family and culture with its own language for disability and support. Those identities affect who is noticed, who is believed, who receives an evaluation, which behavior is punished, and what services are available. \citet{harvey2018school} reinforces that school is not a neutral setting; its routines and expectations reward some bodyminds while pressuring others to mask. \textit{Best Kept Secret} adds the material consequences of race, poverty, geography, and transition access to that analysis \citep{buck2013best}.

This is why lived experience cannot be treated as an optional add-on after the ``real'' academic analysis. Professional knowledge is incomplete without it. Institutional knowledge remains incomplete when it excludes family knowledge. Student behavior is incomplete evidence when adults refuse to consider context. Critical autism studies, historical accounts, genetic discourse, educator guides, family resources, and media narratives all reveal different pieces; they also carry different power to define the person \citep{orsiniDavidson2013critical,nadesan2013genetics,harvey2018story,autismSpeaksReads}.

A culturally responsive understanding therefore requires more than adding identity words to a paragraph. It requires asking whose observations become data, whose language becomes diagnostic, whose story becomes representative, and whose goals shape support.
